% 古籍版式要素总览图的底稿（issue #195）
% 由同目录 annotate.py 编译并加编号标注，生成 anatomy-zh.png / anatomy-en.png，
% 供 wiki「古籍版式要素」页使用。整版不拆筒子页，以便同时看到左右半叶与版心。
% 改动本文件后须重跑 annotate.py，并核对其中 MARKS 的标记位置。
\documentclass[四库全书]{ltc-guji}
\setmainfont{TW-Kai}
\关闭分页
\开启句读模式
\banxinSetup{书耳=true}
\title{欽定四庫全書}
\chapter{史記卷一}

\begin{document}
\begin{正文}
\印章[page=1, opacity=0.75, color=black, xshift=30.9cm, yshift=22.5cm, width=3.2cm]{文渊阁宝印.png}%
\眉批[height=5, x=150pt, y=40pt]{眉批書於天頭可多列}
欽定四庫全書

\begin{段落}[indent=1]
史記卷一
\end{段落}

\begin{段落}[indent=7]
漢太史令司馬遷撰
\end{段落}

黃帝者，少典之子，\夹注{集解徐廣曰號有熊索隠曰少典者諸侯國號}姓公孫，名曰軒轅。生而神靈，弱而能言，幼而徇齊，長而敦敏，成而聰明。軒轅之時，神農氏世衰。諸侯相侵伐，暴虐百姓，而神農氏弗能征。於是軒轅乃習用干戈，以征不享，諸侯咸\侧批{側批在行間}來賓從。而蚩尤最為暴，莫能伐。炎帝欲侵陵諸侯，諸侯咸歸軒轅。軒轅乃修德振兵，治五氣，蓺五種，撫萬民，度四方，教熊羆貔貅貙虎，以與炎帝戰於阪泉之野。三戰，然後得其志。蚩尤作亂，不用帝命。於是黃帝乃徵師諸侯，與蚩尤戰於涿鹿之野，遂禽殺蚩尤。而諸侯咸尊軒轅為天子，代神農氏，是為黃帝。天下有不順者，黃帝從而征之，平者去之，披山通道，未嘗寧居。
\end{正文}
\end{document}
